\clearpage
\oldpage[II]{217}
\markboth{ERRATA AND ADDENDA, LIST 1}{ERRATA AND ADDENDA, LIST 1}
\phantomsection
\addcontentsline{toc}{section}{Errata and Addenda (List 1)}
\section*{Errata and Addenda}
\thispagestyle{plain}
\begin{center}
\large (List 1)
\end{center}

\subsection*{Chapter 0}

\paragraph{(0, 2.1.3).}
Line 15 from the bottom of p.~21, replace ``inter-section'' by
``intersection''.

\paragraph{(0, 2.1.6).}
Line 13 of p.~22, replace
\[
  Z_i\cap C\left(\bigcup_{j\ne i}Z_j\right)
  \quad\hbox{by}\quad
  \bigcup_{j\ne i}Z_j.
\]

\paragraph{(0, 3.1.6).}
Line 17 from the bottom of p.~25, replace
$\Gamma(U,\sh{F})$ by $\Gamma(V,\sh{F})$.

\paragraph{(0, 4.1.1).}
Line 12 of p.~36, replace $\sh{G}\to\sh{F}$ by
$\sh{B}\to\sh{A}$.

\paragraph{(0, 4.2.1).}
Line 11 from the bottom of p.~39, replace ``sheaves'' by ``sheaf''.

\paragraph{(0, 5.1.2).}
After line 14 of p.~45, add:

If $\sh{G}$ is an $\sh{O}_Y$-module generated by its sections over
$Y$, then $f^*(\sh{G})$ is generated by its sections over
$X$, since $f^*$ is a right exact functor.

\paragraph{(0, 5.2.4).}
Line 18 from the bottom of p.~46, replace $f_U$ by $f_U^*$.

\paragraph{(0, 5.3.4).}
Before line 16 from the bottom of p.~47, add:

If
\[
  \sh{A}\to\sh{B}\to\sh{C}\to\sh{D}\to\sh{E}
\]
is an exact sequence of $\sh{O}_X$-modules, and if
$\sh{A},\sh{B},\sh{D},\sh{E}$ are coherent, then $\sh{C}$ is
coherent.

\paragraph{(0, 5.4.1).}
After line 2 of p.~49, add:

Suppose $\sh{O}_X$ coherent, and let $\sh{F}$ be a coherent
$\sh{O}_X$-module. Then, if at a point $x\in X$, $\sh{F}_x$ is a free
$\sh{O}_x$-module of rank $n$, there exists a neighbourhood $U$ of $x$
such that $\sh{F}|U$ is locally free of rank $n$; indeed,
$\sh{F}_x$ is then isomorphic to $\sh{O}_x^n$, and the proposition follows
from (0, 5.2.7).

\paragraph{(0, 6.6.2).}
Replace lines 10 and 11 of p.~59 by:

Indeed, this follows from (0, 6.4.1, d)).

\paragraph{(0, 6.7.1).}
Before line 10 from the bottom of p.~59, add:

If $f$ is a flat morphism, we also say that $X$ is \emph{flat over
$Y$}, or \emph{$Y$-flat}.

\paragraph{(0, 7.2.9).}
Line 11 from the bottom of p.~65, replace $M_n$ by $M_{n-1}$.

\subsection*{Chapter I}

\paragraph{(I, 1.2.1).}
Line 18 of p.~83, replace (twice) $K$ by $k$.

\paragraph{(I, 1.2.7).}
Line 16 from the bottom of p.~84, replace $A$ by $A'$ (twice) and
$\mathfrak N$ by $\mathfrak N'$; line 17 from the bottom, replace $X$ by
$X'$.

\paragraph{(I, 1.7.3) and (I, 2.2.4).}
The statements given in these numbers can (following a remark due to
J.~Tate) be generalized as follows:

\subsection*{1.8. Morphisms of ringed spaces in local rings into affine schemes}
\label{ega2:addendum:subsection:I.1.8}

\begin{proposition}[1.8.1]
\label{ega2:addendum:I.1.8.1}
Let $(S,\sh{O}_S)$ be an affine scheme and $(X,\sh{O}_X)$ a ringed space
in local rings. There is then a canonical bijection from the set
of ring homomorphisms
\oldpage[II]{218}
\markboth{A. GROTHENDIECK --- Chap. II}{A. GROTHENDIECK --- Chap. II}
\[
  \Gamma(S,\sh{O}_S)\to\Gamma(X,\sh{O}_X)
\]
onto the set of morphisms of ringed spaces
$(\psi,\theta):(X,\sh{O}_X)\to(S,\sh{O}_S)$ such that, for every
$x\in X$, $\theta_x^\sharp:\sh{O}_{\psi(x)}\to\sh{O}_x$ is a
local homomorphism.
\end{proposition}

Note first that if $(X,\sh{O}_X)$, $(S,\sh{O}_S)$ are two arbitrary
ringed spaces, a morphism $(\psi,\theta)$ from $(X,\sh{O}_X)$
to $(S,\sh{O}_S)$ canonically defines a ring homomorphism
$\Gamma(\theta):\Gamma(S,\sh{O}_S)\to\Gamma(X,\sh{O}_X)$; whence a
first map
\begin{equation}
\label{ega2:addendum:I.1.8.1.1}
  \rho:\Hom((X,\sh{O}_X),(S,\sh{O}_S))
  \to\Hom(\Gamma(S,\sh{O}_S),\Gamma(X,\sh{O}_X)).
\tag{1.8.1.1}
\end{equation}
Conversely, under the hypotheses of the statement, set
$A=\Gamma(S,\sh{O}_S)$, and consider a ring homomorphism
$\vphi:A\to\Gamma(X,\sh{O}_X)$. For every $x\in X$, it is clear that
the set of $f\in A$ such that $\vphi(f)(x)=0$ is a prime ideal
of $A$, since $\sh{O}_x/\mathfrak m_x=\kres(x)$ is a field; it is
therefore an element of $S=\Spec(A)$, which we will again denote by
${}^a\vphi(x)$. Moreover, for every $f\in A$, we have by definition
(0, 5.5.2)
\[
  {}^a\vphi^{-1}(D(f))=X_f,
\]
which proves that ${}^a\vphi$ is a continuous map $X\to S$.
Next let us define a homomorphism
\[
  \widetilde{\vphi}:\sh{O}_S\to{}^a\vphi_*(\sh{O}_X)
\]
of $\sh{O}_S$-modules; for every $f\in A$, we have
$\Gamma(D(f),\sh{O}_S)=A_f$ (I, 1.3.6); for every $s\in A$, we assign
to $s/f\in A_f$ the element
\[
  (\vphi(s)|X_f)(\vphi(f)|X_f)^{-1}
  \in\Gamma(X_f,\sh{O}_X)
  =\Gamma(D(f),{}^a\vphi_*(\sh{O}_X)),
\]
and we verify immediately (by passing from $D(f)$ to $D(fg)$) that this
indeed defines a homomorphism of $\sh{O}_S$-modules, whence a morphism
$({}^a\vphi,\widetilde{\vphi})$ of ringed spaces. Moreover, with the
same notation, and setting $y={}^a\vphi(x)$ for simplicity, we see
immediately (0, 3.7.1) that we have
\[
  \widetilde{\vphi}_x^\sharp(s_y/f_y)
  =(\vphi(s)_x)(\vphi(f)_x)^{-1};
\]
since the relation $s_y\in\mathfrak m_y$ is by definition equivalent
to $\vphi(s)_x\in\mathfrak m_x$, we see that
$\widetilde{\vphi}_x^\sharp$ is a local homomorphism
$\sh{O}_y\to\sh{O}_x$, and we have thus defined a second map
\[
  \sigma:\Hom(\Gamma(S,\sh{O}_S),\Gamma(X,\sh{O}_X))\to\mathfrak L,
\]
where $\mathfrak L$ is the set of morphisms
$(\psi,\theta):(X,\sh{O}_X)\to(S,\sh{O}_S)$ such that
$\theta_x^\sharp$ is local for every $x\in X$.

It remains to prove that $\sigma$ and $\rho$ (restricted to $\mathfrak L$)
are inverse to each other; now, the definition of
$\widetilde{\vphi}$ shows immediately that
$\Gamma(\widetilde{\vphi})=\vphi$, and consequently $\rho\circ\sigma$ is
the identity. To see that $\sigma\circ\rho$ is the identity, let us start
from a morphism $(\psi,\theta)\in\mathfrak L$ and let
$\vphi=\Gamma(\theta)$; the hypothesis on $\theta_x^\sharp$ allows us to
deduce from this homomorphism, by passage to the quotients, a
monomorphism
\[
  \theta^x:\kres(\psi(x))\to\kres(x)
\]
such that for every section $f\in A=\Gamma(S,\sh{O}_S)$ we have
$\theta^x(f(\psi(x)))=\vphi(f)(x)$; the relation $f(\psi(x))=0$ is therefore
equivalent to $\vphi(f)(x)=0$, which already proves that
${}^a\vphi=\psi$. On the other hand, the definitions imply that the
diagram
\[
  \xymatrix{
    A\ar[r]^\vphi\ar[d] &
    \Gamma(X,\sh{O}_X)\ar[d]\\
    A_{\psi(x)}\ar[r]^{\theta_x^\sharp} &
    \sh{O}_x
  }
\]
is commutative, and the same holds for the analogous diagram in which
$\theta_x^\sharp$ is replaced by
$\widetilde{\vphi}_x^\sharp$, whence
$\widetilde{\vphi}_x^\sharp=\theta_x^\sharp$ (0, 1.2.4), and consequently
$\widetilde{\vphi}=\theta$.

\begin{env}[1.8.2]
\label{ega2:addendum:I.1.8.2}
When $(X,\sh{O}_X)$ and $(Y,\sh{O}_Y)$ are ringed spaces in
local rings, we will have to consider the morphisms
$(\psi,\theta):(X,\sh{O}_X)\to(Y,\sh{O}_Y)$ such that, for every
$x\in X$, $\theta_x^\sharp$ is a local homomorphism
$\sh{O}_{\psi(x)}\to\sh{O}_x$.

When from now on we speak
\oldpage[II]{219}
\markboth{ERRATA AND ADDENDA, LIST 1}{ERRATA AND ADDENDA, LIST 1}
of a \emph{morphism of ringed spaces
in local rings}, it is always such a morphism that will be meant;
with this definition of morphisms, it is clear that the ringed spaces
in local rings form a category; for two objects
$X,Y$ of this category, $\Hom(X,Y)$ will therefore denote the set of
morphisms of ringed spaces in local rings from $X$ to $Y$
(the set denoted $\mathfrak L$ in (1.8.1)); when we have to
consider the set of morphisms of ringed spaces from $X$ to
$Y$, we will denote it by $\Hom_{\mathrm{an}}(X,Y)$ to avoid any
confusion. The map (1.8.1.1) is therefore written
\begin{equation}
\label{ega2:addendum:I.1.8.2.1}
  \rho:\Hom_{\mathrm{an}}(X,Y)
  \to\Hom(\Gamma(Y,\sh{O}_Y),\Gamma(X,\sh{O}_X)),
\tag{1.8.2.1}
\end{equation}
and its restriction
\begin{equation}
\label{ega2:addendum:I.1.8.2.2}
  \rho':\Hom(X,Y)
  \to\Hom(\Gamma(Y,\sh{O}_Y),\Gamma(X,\sh{O}_X))
\tag{1.8.2.2}
\end{equation}
is a map functorial in $X$ and $Y$ in the category of ringed
spaces in local rings.
\end{env}

\begin{corollary}[1.8.3]
\label{ega2:addendum:I.1.8.3}
Let $(Y,\sh{O}_Y)$ be a ringed space in local rings. For $Y$
to be an affine scheme, it is necessary and sufficient that for every
ringed space in local rings $(X,\sh{O}_X)$, the map (1.8.2.2) be
bijective.
\end{corollary}

Proposition (1.8.1) shows that the condition is necessary.
Conversely, if we suppose it satisfied and set
$A=\Gamma(Y,\sh{O}_Y)$, it follows from the hypothesis and from (1.8.1) that
the functors $X\mapsto\Hom(X,Y)$ and
$X\mapsto\Hom(X,\Spec(A))$, from the category of ringed spaces in
local rings to that of sets, are isomorphic. We know that
this implies the existence of a canonical isomorphism
$Y\to\Spec(A)$ (cf.~0, 8).

\begin{env}[1.8.4]
\label{ega2:addendum:I.1.8.4}
Let $S=\Spec(A)$ be an affine scheme; denote by $(S',A')$ the ringed
space whose underlying space consists of a single point and whose
structural sheaf $A'$ is the (necessarily simple) sheaf on $S'$ defined
by the ring $A$. Let $\pi:S\to S'$ be the unique map from $S$ to
$S'$; note on the other hand that for every open subset $U$ of $S$ we have a
canonical map
$\Gamma(S',A')=\Gamma(S,\sh{O}_S)\to\Gamma(U,\sh{O}_S)$, which therefore
defines a $\pi$-morphism $\iota:A'\to\sh{O}_S$ of sheaves of rings.
We have thus canonically defined a morphism of ringed spaces
$i=(\pi,\iota):(S,\sh{O}_S)\to(S',A')$. For every $A$-module $M$, we
denote by $M'$ the simple sheaf on $S'$ defined by $M$, which
is evidently an $A'$-module. It is clear that we have
$i_*(\widetilde M)=M'$ (I, 1.3.7).
\end{env}

\begin{lemma}[1.8.5]
\label{ega2:addendum:I.1.8.5}
With the notation of (1.8.4), for every $A$-module $M$, the
canonical functorial $\sh{O}_S$-homomorphism (0, 4.4.3.3)
\begin{equation}
\label{ega2:addendum:I.1.8.5.1}
  i^*(i_*(\widetilde M))\to\widetilde M
\tag{1.8.5.1}
\end{equation}
is an isomorphism.
\end{lemma}

Indeed, the two sides of (1.8.5.1) are right exact (the
functor $M\mapsto i_*(\widetilde M)$ being evidently exact) and
commute with direct sums; by considering $M$ as the cokernel of a
homomorphism $A^{(I)}\to A^{(J)}$, we are reduced to proving the lemma
in the case $M=A$, and it is evident in that case.

\begin{corollary}[1.8.6]
\label{ega2:addendum:I.1.8.6}
Let $(X,\sh{O}_X)$ be a ringed space and $u:X\to S$ a morphism
of ringed
\oldpage[II]{220}
\markboth{A. GROTHENDIECK --- Chap. II}{A. GROTHENDIECK --- Chap. II}
spaces. For every $A$-module $M$, we have (with the notation of (1.8.4)) a
canonical functorial isomorphism of $\sh{O}_X$-modules
\begin{equation}
\label{ega2:addendum:I.1.8.6.1}
  u^*(\widetilde M)\isoto u^*(i^*(M')).
\tag{1.8.6.1}
\end{equation}
\end{corollary}

\begin{corollary}[1.8.7]
\label{ega2:addendum:I.1.8.7}
Under the hypotheses of (1.8.6), we have, for every $A$-module $M$ and every
$\sh{O}_X$-module $\sh{F}$, a canonical isomorphism, functorial in
$M$ and $\sh{F}$,
\begin{equation}
\label{ega2:addendum:I.1.8.7.1}
  \Hom_{\sh{O}_S}(\widetilde M,u_*(\sh{F}))
  \isoto\Hom_A(M,\Gamma(X,\sh{F})).
\tag{1.8.7.1}
\end{equation}
\end{corollary}

Indeed, by virtue of (0, 4.4.3) and of Lemma (1.8.5), we have a
canonical isomorphism of bifunctors
\[
  \Hom_{\sh{O}_S}(\widetilde M,u_*(\sh{F}))
  \isoto\Hom_{A'}(M',i_*(u_*(\sh{F}))),
\]
and it is clear that the right-hand side is none other than
$\Hom_A(M,\Gamma(X,\sh{F}))$. Note that the canonical homomorphism
(1.8.7.1) assigns to every $\sh{O}_S$-homomorphism
$h:\widetilde M\to u_*(\sh{F})$ (in other words, to every $u$-morphism
$\widetilde M\to\sh{F}$) the $A$-homomorphism
\[
  \Gamma(h):M\to\Gamma(S,u_*(\sh{F}))=\Gamma(X,\sh{F}).
\]

\begin{env}[1.8.8]
\label{ega2:addendum:I.1.8.8}
With the notation of (1.8.4), it is clear (0, 4.1.1) that every
morphism of ringed spaces $(\psi,\theta):X\to S'$ is equivalent to the
data of a ring homomorphism
$A\to\Gamma(X,\sh{O}_X)$. We can therefore also interpret (1.8.1)
as defining a canonical bijection
\[
  \Hom(X,S)\isoto\Hom(X,S')
\]
(where of course the right-hand side consists of morphisms of ringed
spaces, since in general $A$ is not a local ring). More
generally, if $X,Y$ are two ringed spaces in local rings and
if $(Y',A')$ is the ringed space whose underlying space consists of a
single point and whose sheaf of rings $A'$ is the simple sheaf
defined by the ring $\Gamma(Y,\sh{O}_Y)$, we can interpret
(1.8.2.1) as a map
\begin{equation}
\label{ega2:addendum:I.1.8.8.1}
  \rho:\Hom_{\mathrm{an}}(X,Y)\to\Hom(X,Y').
\tag{1.8.8.1}
\end{equation}
The result of (1.8.3) can therefore be interpreted as saying that the affine
schemes are characterized among the ringed spaces in local rings
as those for which the restriction of $\rho$ to $\Hom(X,Y)$:
\begin{equation}
\label{ega2:addendum:I.1.8.8.2}
  \rho':\Hom(X,Y)\to\Hom(X,Y')
\tag{1.8.8.2}
\end{equation}
is bijective for every ringed space in local rings $X$. In a
later chapter, we will generalize this definition, which will
make it possible to associate to an arbitrary ringed space $Z$ (and no longer
only to a ringed space whose underlying space consists of a single
point) a ringed space in local rings which we will again call
$\Spec(Z)$; this will be the starting point of a ``relative'' theory of
preschemes over arbitrary ringed spaces, extending the
results of Chapter~I.
\end{env}

\begin{env}[1.8.9]
\label{ega2:addendum:I.1.8.9}
We can consider the pairs $(X,\sh{F})$ consisting of a ringed space
in local rings $X$ and an $\sh{O}_X$-module $\sh{F}$ as forming
a category, a morphism $(X,\sh{F})\to(Y,\sh{G})$ of this category
being a pair $(u,h)$ consisting of a morphism of ringed
\oldpage[II]{221}
\markboth{ERRATA AND ADDENDA, LIST 1}{ERRATA AND ADDENDA, LIST 1}
spaces in local rings $u:X\to Y$ and a $u$-morphism
$h:\sh{G}\to\sh{F}$ of modules; these
morphisms (for fixed $(X,\sh{F})$ and $(Y,\sh{G})$) form a set
which we will denote by $\Hom((X,\sh{F}),(Y,\sh{G}))$; the map
$(u,h)\mapsto(\rho'(u),\Gamma(h))$ is a canonical map
\begin{equation}
\label{ega2:addendum:I.1.8.9.1}
\begin{split}
  \Hom((X,\sh{F}),(Y,\sh{G}))\to
  \Hom\bigl(
    &(\Gamma(Y,\sh{O}_Y),\Gamma(Y,\sh{G})),\\
    &(\Gamma(X,\sh{O}_X),\Gamma(X,\sh{F}))
  \bigr),
\end{split}
\tag{1.8.9.1}
\end{equation}
functorial in $(X,\sh{F})$ and $(Y,\sh{G})$, the right-hand side being
the set of di-homomorphisms corresponding to the rings and
modules under consideration (0, 1.0.2).
\end{env}

\begin{corollary}[1.8.10]
\label{ega2:addendum:I.1.8.10}
Let $Y$ be a ringed space in local rings and $\sh{G}$ an
$\sh{O}_Y$-module. For $Y$ to be an affine scheme and $\sh{G}$ a
quasi-coherent $\sh{O}_Y$-module, it is necessary and sufficient that for
every pair $(X,\sh{F})$ consisting of a ringed space in local rings $X$
and an $\sh{O}_X$-module $\sh{F}$, the canonical map (1.8.9.1) be
bijective.
\end{corollary}

We leave it to the reader to develop the argument, modelled
on that of (1.8.3), in which one uses (1.8.1) and (1.8.7).

\begin{remark}[1.8.11]
\label{ega2:addendum:I.1.8.11}
The statements (I, 1.7.3), (I, 1.7.4) and (I, 2.2.4) are particular
cases of (1.8.1), as is the definition (I, 1.6.1); likewise
(I, 2.2.5) follows from (1.8.7). Proposition (1.8.7) also implies
(I, 1.6.3) (and consequently (I, 1.6.4)) as a particular case, since if $X$
is affine and $\Gamma(X,\sh{F})=N$, the functors
\[
  M\longmapsto
  \Hom_{\sh{O}_S}(\widetilde M,u_*(\widetilde N))
  \quad\hbox{and}\quad
  M\longmapsto
  \Hom_{\sh{O}_S}(\widetilde M,(N_{[\vphi]})\supertilde)
\]
(where $\vphi:A\to\Gamma(X,\sh{O}_X)$ corresponds to $u$) are isomorphic
by (1.8.7) and (I, 1.3.8). Finally, (I, 1.6.5) (and consequently (I, 1.6.6))
follows from (1.8.6) and from the fact that for every $f\in A$, the
$A_f$-modules
\[
  N'\otimes_{A'}A_f
  \quad\hbox{and}\quad
  (N'\otimes_{A'}A)_f
\]
(notation of (I, 1.6.5)) are canonically isomorphic.
\end{remark}

\paragraph{Insert after (I, 3.2.8):}

\begin{env}[3.2.9]
\label{ega2:addendum:I.3.2.9}
It follows from (1.8.1) that (I, 3.2.2) can be made more precise in the
following way: $Z=\Spec(B\otimes_A C)$ is a product of
$X=\Spec(B)$ and $Y=\Spec(C)$ not only in the category of $S$-preschemes,
but also in the category of ringed spaces in local rings
over $S$ (with a definition of $S$-morphisms modelled on
that of (I, 2.5.2)). The proof of (I, 3.2.6) then in fact proves
that for two arbitrary $S$-preschemes $X,Y$, the prescheme
$X\times_S Y$ is the product of $X$ and $Y$ not only in the
category of $S$-preschemes, but also in the category of ringed
spaces in local rings over the prescheme $S$.
\end{env}

\paragraph{(I, 4.4.5).}
Line 14 from the bottom of p.~126, replace $B$ by $A$; line 13 from the
bottom, replace ``$A$-algebra'' by ``$B$-algebra''.

\paragraph{(I, 5.3.5).}
Line 2 from the bottom of p.~132, replace $Y\to S$ by $g:Y\to S$.

\paragraph{(I, 6.3.2.1).}
Line 6 from the bottom of p.~144, read:
\[
  D(g_i)\subset W.
\]

\paragraph{(I, 7.3.8).}
Replace the present text, which is erroneous, by the following:

Let $X,Y$ be two integral preschemes, which implies that
$\sh{R}(X)$ (resp. $\sh{R}(Y)$) is a quasi-coherent $\sh{O}_X$-module
(resp. a quasi-coherent $\sh{O}_Y$-module) (I, 7.3.3). Let $f:X\to Y$ be a
dominant morphism; then there exists a canonical homomorphism of
$\sh{O}_X$-modules
\oldpage[II]{222}
\markboth{A. GROTHENDIECK --- Chap. II}{A. GROTHENDIECK --- Chap. II}
\begin{equation}
\label{ega2:addendum:I.7.3.8.1}
  \tau:f^*(\sh{R}(Y))\to\sh{R}(X).
\tag{7.3.8.1}
\end{equation}

Suppose first that $X=\Spec(A)$ and $Y=\Spec(B)$ are affine, with
integral rings $A$ and $B$, so that $f$ corresponds to an injective
homomorphism $B\to A$, which extends to a monomorphism $L\to K$ from the
field of fractions $L$ of $B$ into the field of fractions $K$ of $A$.
The homomorphism (7.3.8.1) then corresponds to the canonical homomorphism
\[
  L\otimes_B A\to K
\]
(I, 1.6.5). In the general case, for every pair of non-empty affine open
subsets $U\subset X$, $V\subset Y$ such that $f(U)\subset V$, we define
in the preceding way a homomorphism $\tau_{U,V}$, and we observe
immediately that if $U'\subset U$, $V'\subset V$,
$f(U')\subset V'$, then $\tau_{U,V}$ extends $\tau_{U',V'}$, whence our
assertion. If $x,y$ are the generic points of $X$ and $Y$
respectively, we have $f(x)=y$,
\[
  \bigl(f^*(\sh{R}(Y))\bigr)_x
  =\sh{O}_y\otimes_{\sh{O}_y}\sh{O}_x
  =\sh{O}_x
\]
(0, 4.3.1), and $\tau_x$ is therefore an isomorphism.

\paragraph{(I, 9.5.2).}
In the last 3 lines of p.~176 and the first 5 lines of
p.~177, replace $X'$ by $Y'$ and $X''$ by $Y''$ throughout.

\paragraph{(I, 10.5.4).}
Line 11 from the bottom of p.~187, replace (10.3.4) by (10.3.5).

\paragraph{(I, 10.6.3).}
Line 3 from the bottom of p.~189, replace $X$ by $\mathfrak X$.

\paragraph{(I, 10.14.2).}
Line 5 of p.~210, replace ``coherent $\sh{O}_{\mathfrak X}$-module'' by
``coherent ideal of $\sh{O}_{\mathfrak X}$''; line 7, replace
``coherent $\sh{O}_{\mathfrak X}$-modules'' by ``coherent ideals of
$\sh{O}_{\mathfrak X}$''.
